% =============================================================================
%  Inference_DayaBay.tex
%  Real Multi-Reactor theta13/Delta m^2 Inference at Daya Bay with tpeanuts
%  Fast Simulation of Neutrino Oscillations in Matter (MSc Thesis)
% =============================================================================
\documentclass[11pt,a4paper,openany]{report}

\usepackage[a4paper,margin=2.6cm,marginparwidth=2.1cm,marginparsep=4mm]{geometry}
\usepackage{tpeanuts-inference}

\addbibresource{references.bib}

\title{\Huge Real Multi-Reactor $\theta_{13}$/$\Delta m^2$ Inference at Daya Bay with \texttt{tpeanuts}\\[4pt]
  \Large Fast Simulation of Neutrino Oscillations in Matter}
\author{Juan Ramon Diaz Santos}
\date{August 2026 --- living document, generated from the source code}

\begin{document}

\begin{titlepage}
  \centering
  \vspace*{1.6cm}
  {\Huge\bfseries\color{dayabayColor} Real Multi-Reactor $\theta_{13}$/$\Delta m^2$\\[4pt] Inference at Daya Bay with \texttt{tpeanuts}\par}
  \vspace{0.8cm}
  {\Large Fast Simulation of Neutrino Oscillations in Matter\par}
  \vspace{1.5cm}
  {\large MSc Thesis (Trabajo Fin de M\'aster) --- Inference Companion Document\par}
  \vspace{0.3cm}
  {\large Juan Ramon Diaz Santos \\ \texttt{diazjuan@alumni.uv.es}\par}
  \vspace{0.6cm}
  {\normalsize Supervisors: Roberto Ruiz de Austri Bazan \quad\&\quad Michele Lucente\par}
  \vfill
  {\normalsize September 2026\par}
  \vspace{0.5cm}
  \begin{tcolorbox}[colback=dayabayColor!5,colframe=dayabayColor,width=0.85\textwidth,sharp corners=all,rounded corners=all,arc=1mm]
    \small
    This document accompanies \texttt{notebooks/inference/inference4\_dayabay.ipynb} and the
    \texttt{tpeanuts.detector.dayabay}/\texttt{tpeanuts.inference} code packages: it documents,
    end to end, the real Daya Bay Collaboration reactor-antineutrino experiment, the official
    public data release it is fitted against, the detector-response and reactor-flux physics
    implemented in \texttt{tpeanuts}, the inference machinery (likelihoods, fits, likelihood
    scans), and the resulting real $\theta_{13}$/$\Delta m^2_{31}$ measurement, compared against
    Daya Bay's own published values and against NuFIT. Blue boxes (\textbf{Theoretical
    Foundation}) contain the physics; green boxes (\textbf{Implementation}) link that physics to
    actual \texttt{tpeanuts} code; orange boxes collect approximations and limitations; terracotta
    boxes (\textbf{Result}) report a concrete numerical outcome of the notebook. Text highlighted
    in \keyresult{yellow} marks key results/equations.
  \end{tcolorbox}
\end{titlepage}

\pagenumbering{roman}
\tableofcontents
\listoftodos[Pending sections and notes]
\clearpage
\pagenumbering{arabic}

% =============================================================================
\chapter{Introduction to the Daya Bay Experiment}
\label{ch:introduction}

\section{Purpose and how to read this document}

This document is the inference-side companion to the \emph{Theoretical
Foundations of \texttt{tpeanuts}} document: where that document develops the
generic neutrino-propagation physics implemented in \texttt{core/} and
\texttt{medium/}, this one documents a single, complete, real-data
application built on top of it -- fitting the reactor mixing angle
$\theta_{13}$ and the mass-squared splitting $\Delta m^2_{31}$ to the real
public data release of the Daya Bay reactor antineutrino experiment.
Every number quoted here that is not explicitly attributed to a fitted
\texttt{tpeanuts} result is a real, published, externally verifiable
quantity (a detector specification, a reactor parameter, or a Daya Bay
Collaboration measurement), traceable through the References chapter.

The five chapters are organised as follows: this chapter introduces the
experiment itself -- geometry, detector design, and detection principle --
with no reference yet to \texttt{tpeanuts}; \cref{ch:data-release}
catalogues the real official data release the code consumes, table by
table; \cref{ch:detector-implementation} is the physics-and-mathematics
core, developing every formula \texttt{tpeanuts.detector.dayabay}
implements; \cref{ch:inference-process} documents the statistical inference
machinery (likelihoods, fits, likelihood scans) and the three fit
strategies used; \cref{ch:results} presents, interprets, and summarises the
numerical results obtained from \texttt{inference4\_dayabay.ipynb},
compared against Daya Bay's own published values and against NuFIT.

\section{Overview}

The Daya Bay Reactor Neutrino Experiment is a multi-detector, multi-reactor
survey of electron-antineutrino disappearance, located at the Daya Bay
nuclear power complex in Shenzhen, Guangdong province, China. It was
designed with a single, focused physics goal: to measure the smallest of
the three PMNS mixing angles, $\theta_{13}$, via the disappearance channel
\begin{equation}
  \bar\nu_e \to \bar\nu_e, \qquad
  P_{\bar\nu_e\bar\nu_e}(E,L) \simeq 1 - \sin^2 2\theta_{13}\,
    \sin^2\!\left(\frac{\Delta m^2_{ee}L}{4E}\right)
    + \left(\theta_{12},\,\Delta m^2_{21}\ \text{terms}\right),
  \label{eq:intro-survival-schematic}
\end{equation}
using only reactor-produced $\bar\nu_e$ (no beam, no accelerator) and pairs
of identically designed detectors at different baselines. Data taking began
in December 2011 and continued through 2020, accumulating more than 5.55
million inverse-beta-decay (IBD) candidates with neutron capture on
gadolinium (nGd) over 3158 days of operation
\parencite{An2017DayaBayFluxSpectrum}. The experiment reported the first
$>5\sigma$ observation of a non-zero $\theta_{13}$ in March 2012
\parencite{An2012DayaBayDiscovery}, and subsequently refined the
measurement to few-percent precision on both $\sin^2 2\theta_{13}$ and the
effective splitting $\Delta m^2_{ee}$ \parencite{An2023DayaBayPrecision}.
This document, and the \texttt{tpeanuts} code it accompanies, is built
against the Collaboration's own full public data release
\parencite{DayaBayZenodo2025}, the same underlying dataset behind
\cite{An2023DayaBayPrecision}.

\begin{notebox}
\textbf{Scope of this document.} \texttt{tpeanuts} reproduces the real
6-reactor/8-detector geometry, the real reactor antineutrino flux (with its
real non-equilibrium and spent-fuel corrections), the real order-1/M
inverse-beta-decay cross section, and the real detector energy response
(IAV, LSNL, Gaussian resolution) at their nominal/central values. It does
\emph{not} reproduce the full correlated systematic covariance matrix of
the official Daya Bay combined analysis, nor the 6AD/7AD sub-periods (only
the 8AD period, all 8 detectors online, is modelled). \Cref{ch:results}
reports, honestly and without adjustment, where this simplification does
and does not matter for the final $\theta_{13}$/$\Delta m^2_{31}$ result.
\end{notebox}

\section{The reactor complex: six real cores, three nuclear power plants}
\label{sec:reactor-geometry}

\begin{theorybox}
The antineutrino sources are six pressurised-water reactor cores, grouped
in pairs at three nuclear power plants along the Daya Bay/Ling Ao
site \parencite{An2017DayaBayFluxSpectrum,DayaBayZenodo2025}:
\begin{itemize}
  \item \textbf{Daya Bay NPP}: reactor cores D1, D2.
  \item \textbf{Ling Ao NPP}: reactor cores L1, L2.
  \item \textbf{Ling Ao-II NPP}: reactor cores L3, L4.
\end{itemize}
In the official data release and throughout this document (and the
\texttt{tpeanuts} code), the six cores are labelled generically
$\mathrm{R1}$--$\mathrm{R6}$. Each core has a real nominal thermal power of
\begin{equation}
  P_{\rm th} \approx 2.895\ \mathrm{GW_{th}},
\end{equation}
a single shared, time-averaged value in the data release used identically
for all 6 cores (the real per-core power history varies week to week; see
\cref{sec:weekly-rate-table}). Antineutrinos are produced by the beta decays
of neutron-rich fission fragments of four fissile isotopes present in the
reactor fuel, $^{235}$U, $^{238}$U, $^{239}$Pu, $^{241}$Pu, with real
time-averaged fission fractions (shared by all 6 cores in this data tier)
\begin{equation}
  f_{^{235}{\rm U}} \approx 0.563,\quad
  f_{^{238}{\rm U}} \approx 0.076,\quad
  f_{^{239}{\rm Pu}} \approx 0.305,\quad
  f_{^{241}{\rm Pu}} \approx 0.056.
\end{equation}
On average, about six antineutrinos are released per fission, but only
those above the inverse-beta-decay threshold ($E_\nu \gtrsim 1.8$~MeV,
\cref{ch:detector-implementation}) are detectable -- roughly two per
fission, integrated over the real tabulated Huber-Mueller spectra
(\cref{ch:data-release}).
\end{theorybox}

\begin{notebox}
Every one of the 6 real cores is treated as a point source located at its
own real, published position; the finite physical extent of a reactor core
introduces a negligible baseline correction (a few centimetres against
baselines of hundreds of metres) and is not modelled separately -- this
matches the official analysis's own documented treatment
\parencite{An2017DayaBayFluxSpectrum}.
\end{notebox}

\section{The detector array: three experimental halls, eight antineutrino detectors}
\label{sec:detector-geometry}

\begin{theorybox}
Eight functionally identical antineutrino detectors (ADs) are deployed
underground in three experimental halls (EH), chosen to sit at baselines
spanning the full oscillatory pattern of \labelcref{eq:intro-survival-schematic}:
\begin{itemize}
  \item \textbf{EH1} (near hall): 2 ADs (AD11, AD12), real flux-weighted
    baseline $\approx 560$~m.
  \item \textbf{EH2} (near hall): 2 ADs (AD21, AD22), real flux-weighted
    baseline $\approx 600$~m.
  \item \textbf{EH3} (far hall): 4 ADs (AD31, AD32, AD33, AD34), real
    flux-weighted baseline $\approx 1640$~m.
\end{itemize}
``Flux-weighted baseline'' means the effective distance obtained by
weighting each of the 6 real (detector, reactor) baselines by that
reactor's real relative antineutrino output -- since every detector
actually sees a superposition of all 6 cores, not a single one
(\cref{eq:dayabay-flux-sum} in \cref{ch:detector-implementation}). EH1 and
EH2 together form the ``near'' halls (only weakly oscillated, used as the
reference spectrum in a near/far comparison, \cref{ch:inference-process});
EH3 is the ``far'' hall, sitting close to the first oscillation minimum for
the bulk of the reactor antineutrino energy spectrum and hence carrying
most of the real $\theta_{13}$-driven deficit.

The detector-naming convention used throughout this document and the
\texttt{tpeanuts} code, \code{AD}$hd$, encodes the hall number $h$ and a
sequential detector index $d$ within that hall directly in the name (e.g.\
\code{AD31} is the first AD installed at EH3) -- this is also how
\texttt{tpeanuts.detector.dayabay.parameters.NEAR\_DETECTORS}/
\texttt{FAR\_DETECTORS} group detectors by hall
(\cref{ch:detector-implementation}).
\end{theorybox}

\subsection*{Detector construction}

\begin{theorybox}
Each AD is built from three nested cylindrical vessels, designed to be as
close to functionally identical across all 8 detectors as physically
achievable (a deliberate design choice: systematic detector-to-detector
differences would otherwise be indistinguishable from a real oscillation
effect in a near/far comparison) \parencite{An2017DayaBayFluxSpectrum}:
\begin{enumerate}
  \item \textbf{Inner Acrylic Vessel (IAV)}, wall thickness $11$~mm,
    containing $0.1\%$ gadolinium-doped liquid scintillator (GdLS) --
    the actual neutrino target volume. Gadolinium's huge thermal-neutron
    capture cross section ($\sigma\sim 4.9\times10^4$~b for natural Gd,
    versus $\sim0.3$~b for hydrogen) is what makes the delayed-coincidence
    tag efficient and fast (\cref{sec:ibd-principle}).
  \item \textbf{Outer Acrylic Vessel (OAV)}, filled with undoped liquid
    scintillator (LS) -- a ``gamma catcher'' that increases the detection
    efficiency for gamma rays (from both the prompt positron-annihilation
    signal and the delayed neutron-capture signal) that escape the smaller
    IAV volume without being fully absorbed there.
  \item \textbf{Outermost stainless-steel tank}, filled with mineral oil,
    hosting a total of $192$ 8-inch photomultiplier tubes (PMTs) radially
    positioned around the two inner vessels, plus specular reflectors above
    and below the outer acrylic vessel to improve light collection
    uniformity.
\end{enumerate}
Three automated calibration units (ACUs), each capable of deploying
radioactive sources along one of three vertical axes, sit atop each AD's
outer tank, providing the real energy-scale/position calibration data
underlying the LSNL and IAV response corrections
(\cref{sec:iav-lsnl-theory}).

Each AD is submerged in a two-zone water-Cherenkov muon-veto system (inner
and outer water shield, IWS/OWS), instrumented with its own PMTs, providing
the real muon-veto efficiency folded into the real daily
\code{eff\_livetime} exposure (\cref{sec:efficiency-exposure}).
\end{theorybox}

\subsection*{The inverse-beta-decay detection principle}
\label{sec:ibd-principle}

\begin{theorybox}
Every real event this experiment (and this document's inference) is built
from is an inverse-beta-decay (IBD) interaction,
\begin{equation}
  \keyresult{\bar\nu_e + p \rightarrow e^+ + n,}
  \label{eq:intro-ibd-reaction}
\end{equation}
occurring on a free (unbound hydrogen) proton in the GdLS target. IBD is
selected via a real, two-part delayed-coincidence signature, occurring
within the same detector:
\begin{description}
  \item[Prompt signal.] The positron deposits its kinetic energy locally and
    then annihilates with an atomic electron, $e^+ + e^- \to 2\gamma$
    ($2\times 511$~keV), giving a real visible (``prompt'') energy
    \begin{equation}
      E_{\rm prompt} \approx T_{e^+} + 2 m_e = E_\nu - \Delta_{np} + m_e,
      \qquad \Delta_{np} = M_n - M_p,
    \end{equation}
    at leading order (the full order-1/M kinematics are developed in
    \cref{ch:detector-implementation}).
  \item[Delayed signal.] The IBD neutron thermalises (over microseconds)
    and is captured, predominantly on a real gadolinium nucleus
    (chosen for this purpose precisely because of its large capture cross
    section), releasing a real $\sim 8$~MeV gamma-ray cascade $\sim 30\
    \mu\mathrm{s}$ later on average -- well separated in both time and
    energy from the prompt signal and from most backgrounds.
  \item[Selection cuts.] A sequence of real cuts (flasher rejection,
    capture-time window, prompt/delayed energy windows, muon veto,
    multiplicity) isolates the real IBD candidate sample; \cref{ch:results}
    reports how tpeanuts' own real IBD-selection efficiency compares to the
    official one.
\end{description}
This delayed-coincidence tag is what makes IBD, among the handful of
reactions a reactor antineutrino can undergo in a scintillator detector,
overwhelmingly the cleanest and most background-suppressible channel --
exactly why every reactor-antineutrino oscillation experiment to date
(Daya Bay, RENO, Double Chooz, KamLAND) is built around it.
\end{theorybox}

\begin{notebox}
The real coincidence window, multiplicity cuts, and calibration procedures
that turn raw PMT waveforms into a validated IBD candidate list are
themselves a substantial piece of the Collaboration's own analysis
infrastructure; \texttt{tpeanuts} does not reproduce this reconstruction
chain, and instead consumes its real, already-processed output directly
(the official released event-count tables, \cref{ch:data-release}) as its
external, unmodified real-data input.
\end{notebox}

\section{Physics goals and data-taking timeline}

\begin{theorybox}
Real data taking proceeded in three real sub-periods, distinguished by how
many of the 8 ADs were physically installed and operating:
\begin{description}
  \item[6AD period] (December 2011 -- July 2012): 6 ADs online (2 at EH1,
    1 at EH2, 3 at EH3) -- the dataset behind the original 2012 discovery
    \parencite{An2012DayaBayDiscovery}.
  \item[8AD period] (October 2012 onward): all 8 ADs online (2 at each near
    hall, 4 at the far hall) -- the largest and most-used real dataset, and
    the \emph{only} period this document's \texttt{tpeanuts} model fits
    (\cref{ch:detector-implementation}).
  \item[7AD period] (a later real sub-interval): AD11 temporarily removed
    from physical data taking.
\end{description}
The physics goal evolved with the data: the 6AD-period near/far rate
comparison established $\theta_{13}\neq0$
\parencite{An2012DayaBayDiscovery}; with the accumulated 8AD statistics, a
full spectral-shape analysis (comparing the real reconstructed
prompt-energy spectrum bin by bin, not just the total rate) additionally
constrains the mass-squared splitting via the energy-dependent position of
the oscillation minimum, reported as the effective splitting $\Delta
m^2_{ee}$ \parencite{An2023DayaBayPrecision}. \Cref{ch:inference-process}
documents both strategies (rate-only and full-spectral-shape) as
implemented in \texttt{tpeanuts}, plus a real near/far \emph{ratio} fit
that recovers both parameters jointly while remaining robust to the
detector-response simplifications made in this project
(\cref{ch:results}).
\end{theorybox}


% =============================================================================
\chapter{The Official Data Release: Tables and Publications}
\label{ch:data-release}

\section{Publications}

\begin{refbox}
\begin{itemize}
  \item \textbf{Data release.} Daya Bay Collaboration, \emph{Full Data
    Release of the Daya Bay Reactor Neutrino Experiment}, v1.0.0, Zenodo,
    DOI:10.5281/zenodo.17587229 (2025) \parencite{DayaBayZenodo2025}. The
    analysis-dataset (TSV) tier is fetched from the official GitHub mirror
    \code{github.com/dayabay-experiment/dayabay-data-official} by
    \code{notebooks/external/reactor/dayabay/DayaBay1\_generator.ipynb}
    (\code{zenodo.org} itself returns HTTP 403 from this project's
    development sandbox) into \code{data/detector/dayabay/}, in the tidy
    \code{.csv} form \texttt{tpeanuts.detector.dayabay} consumes.
  \item \textbf{Main precision measurement.} F.~P.~An et al.\ (Daya Bay
    Collaboration), \emph{Precision Measurement of Reactor Antineutrino
    Oscillation at Kilometer-Scale Baselines by Daya Bay}, Phys.\ Rev.\
    Lett.\ \textbf{130}, 161802 (2023) \parencite{An2023DayaBayPrecision}.
    Publishes the $\sin^2 2\theta_{13}$/$\Delta m^2_{ee}$ best fit compared
    against this document's own inference in \cref{ch:results}.
  \item \textbf{Discovery paper.} F.~P.~An et al.\ (Daya Bay
    Collaboration), \emph{Observation of Electron-Antineutrino
    Disappearance at Daya Bay}, Phys.\ Rev.\ Lett.\ \textbf{108}, 171803
    (2012) \parencite{An2012DayaBayDiscovery}. Source of the original
    near/far rate-comparison strategy \cref{ch:inference-process}
    reproduces.
  \item \textbf{Flux and spectrum measurement.} F.~P.~An et al.\ (Daya Bay
    Collaboration), \emph{Improved measurement of the reactor antineutrino
    flux and spectrum at Daya Bay}, Chin.\ Phys.\ C \textbf{41}, 013002
    (2017) \parencite{An2017DayaBayFluxSpectrum}. Source of an
    independently published per-detector IBD rate table (its Table 5) used
    in \cref{ch:results} as an external cross-check of this project's
    absolute-normalization audit, and of the Reactor Antineutrino Anomaly
    measurement discussed in \cref{ch:detector-implementation}.
  \item \textbf{External comparison: NuFIT.} I.~Esteban et al., \emph{The
    fate of hints: updated global analysis of three-flavor neutrino
    oscillations}, JHEP \textbf{12}, 216 (2024), arXiv:2410.05380
    \parencite{EstebanNuFit2024} -- base analysis publication for NuFIT 6.1
    (2025), \url{www.nu-fit.org}. Source of the fixed solar parameters used
    throughout \cref{ch:detector-implementation}, and of the independent
    NuFIT $\theta_{13}$ comparison in \cref{ch:results}.
\end{itemize}
\end{refbox}

\section{Data-release structure and format conventions}

\begin{theorybox}
The official release ships four parallel data-format tiers (\code{npz},
\code{hdf5}, \code{tsv.bz2}, \code{root}), each containing the identical
physical content; \texttt{tpeanuts} consumes the \code{tsv.bz2} tier only,
since it requires no format-specific reader beyond the standard-library
\code{bz2} module and \code{pandas}. Three real naming conventions appear
in the tables catalogued below:
\begin{itemize}
  \item \code{\{filename\}\_\{key\}.tsv.bz2} -- a single array, named by
    its key (e.g.\ \code{detector\_iav\_matrix\_iav\_matrix.tsv.bz2}).
  \item \code{\{filename\}.tsv/\{key\}.tsv.bz2} -- a folder of related
    arrays, one file per key (e.g.\ per-detector IBD spectra).
  \item \code{\{filename\}.tsv/\{filename\}\_\{key\}.tsv.bz2} -- the same,
    with the folder name repeated in each file for extra clarity when
    nested paths are otherwise ambiguous.
\end{itemize}
Real parameter files (\code{parameters/*.yaml}) additionally carry a
\code{format} field describing their own uncertainty convention
(\code{value}: point value only; \code{sigma\_absolute}: value $\pm$
absolute 1$\sigma$; \code{sigma\_percent}: value $\pm$ percent 1$\sigma$)
and a \code{state} field (\code{fixed} or \code{variable}, i.e.\ whether the
official combined analysis treats that parameter as a free nuisance).
\end{theorybox}

\section{Table catalogue}

\Cref{tab:dayabay-tables-geometry,tab:dayabay-tables-response,tab:dayabay-tables-observed,tab:dayabay-tables-truth}
list every real table this project fetches and the \texttt{tpeanuts}
loader that reads it (all in \texttt{tpeanuts.detector.dayabay.io}),
grouped by category; the physics each table encodes is developed in full
in \cref{ch:detector-implementation}.

\newcommand{\dbcell}[1]{\parbox[t]{0.47\textwidth}{\raggedright #1}}

\begin{table}[htb]
\centering
\caption{Real geometry and reactor-flux tables.}
\label{tab:dayabay-tables-geometry}
\begin{tabular}{@{}l l l@{}}
\toprule
\textbf{Cached file} & \textbf{Real contents} & \textbf{Loader} \\
\midrule
\code{baselines.csv} & \dbcell{(detector, reactor) baseline in km; real $8\times6=48$-row geometry matrix.} & \pyfunc{load\_baselines} \\[4pt]
\code{reactor\_parameters.csv} & \dbcell{Real fission fractions, energy released per fission (4 isotopes), nominal thermal power $2.895$~GW.} & \pyfunc{load\_reactor\_parameters} \\[4pt]
\code{huber\_mueller\_spectra.csv} & \dbcell{Real per-isotope antineutrino spectra \parencite{Huber2011,Huber2012Erratum,Mueller2011}, $50$~keV mesh, $1.75$--$13.0$~MeV.} & \pyfunc{load\_huber\_mueller\_spectra} \\[4pt]
\code{nonequilibrium\_correction.csv} & \dbcell{Real per-isotope non-equilibrium flux correction ($^{235}$U, $^{239}$Pu, $^{241}$Pu only).} & \pyfunc{load\_nonequilibrium\_correction} \\[4pt]
\code{snf\_correction.csv} & \dbcell{Real per-reactor spent-nuclear-fuel flux correction, all 6 cores.} & \pyfunc{load\_snf\_correction} \\[4pt]
\code{neutrino\_rate\_weekly.csv} & \dbcell{Real weekly-resolved per-reactor antineutrino rate, full 2011--2020 data-taking period.} & \pyfunc{load\_neutrino\_rate\_weekly} \\
\bottomrule
\end{tabular}
\end{table}

\begin{table}[htb]
\centering
\caption{Real detector-response and IBD-cross-section tables.}
\label{tab:dayabay-tables-response}
\begin{tabular}{@{}l l l@{}}
\toprule
\textbf{Cached file} & \textbf{Real contents} & \textbf{Loader} \\
\midrule
\code{n\_protons.csv} & \dbcell{Real target free-proton count per detector ($\approx 1.43\times10^{30}$ nominal, $\pm$ small per-detector correction).} & \pyfunc{load\_n\_protons} \\[4pt]
\code{eres\_parameters.csv} & \dbcell{Real 3-term Gaussian energy-resolution coefficients $(a,b,c)$.} & \pyfunc{load\_eres\_parameters} \\[4pt]
\code{detector\_efficiency.txt} & \dbcell{Real flat IBD-selection efficiency, $\varepsilon\approx0.8025$.} & \pyfunc{load\_detector\_efficiency} \\[4pt]
\code{ibd\_constants.csv} & \dbcell{Real IBD cross-section constants: $f=1.0$, $g=1.2701$, $f_2=3.706$, PhaseSpaceFactor$=1.71465$.} & \pyfunc{load\_ibd\_constants} \\[4pt]
\code{iav\_matrix.csv} & \dbcell{Real $240\times240$ IAV energy-redistribution matrix, 0.05~MeV bins, 0--12~MeV, columns summing to 1.} & \pyfunc{load\_iav\_matrix} \\[4pt]
\code{lsnl\_curve\_nominal.csv} & \dbcell{Real LSNL nominal nonlinearity curve $f(E)=E_{\rm reco}/E_{\rm true}$.} & \pyfunc{load\_lsnl\_curve} \\[4pt]
\code{lsnl\_curve\_pulls.csv} & \dbcell{Real 4 LSNL systematic pull curves $f_k(E)$, $k=0..3$.} & \pyfunc{load\_lsnl\_curve\_pulls} \\
\bottomrule
\end{tabular}
\end{table}

\begin{table}[htb]
\centering
\caption{Real observed-data and background tables (8AD period).}
\label{tab:dayabay-tables-observed}
\begin{tabular}{@{}l l l@{}}
\toprule
\textbf{Cached file} & \textbf{Real contents} & \textbf{Loader} \\
\midrule
\code{ibd\_spectra/ibd\_spectrum\_\{AD\}.csv} & \dbcell{Real observed IBD prompt-energy candidate spectra, 0.05~MeV bins, 0--12~MeV, one file per detector.} & \pyfunc{load\_ibd\_spectrum} \\[4pt]
\code{background\_rates\_8AD.csv} & \dbcell{Real background rates + uncertainties, 5 categories $\times$ 8 detectors.} & \pyfunc{load\_background\_rates} \\[4pt]
\code{background\_shapes/spectrum\_shape\_\{cat\}\_\{AD\}.csv} & \dbcell{Real normalised background spectral shapes (integrate to 1), same binning as the IBD spectra.} & \pyfunc{load\_background\_shape} \\[4pt]
\code{exposure\_8AD.csv} & \dbcell{Real per-detector effective livetime (\code{livetime}$\times$\code{eff}), summed over 8AD-period days.} & \pyfunc{load\_exposure} \\[4pt]
\code{final\_erec\_bin\_edges.csv} & \dbcell{Real (non-uniform) analysis binning, 0.7--12.0~MeV, 27 edges.} & \pyfunc{load\_final\_erec\_bin\_edges} \\
\bottomrule
\end{tabular}
\end{table}

\begin{table}[htb]
\centering
\caption{Real published-truth and normalization tables.}
\label{tab:dayabay-tables-truth}
\begin{tabular}{@{}l l l@{}}
\toprule
\textbf{Cached file} & \textbf{Real contents} & \textbf{Loader} \\
\midrule
\code{survival\_probability\_truth.csv} & \dbcell{Daya Bay's own real published best fit ($\sin^2 2\theta_{13}$, $\Delta m^2_{32}$) and solar input ($\sin^2 2\theta_{12}$, $\Delta m^2_{21}$), for comparison only.} & \pyfunc{load\_survival\_probability\_truth} \\[4pt]
\code{global\_normalization.txt} & \dbcell{Daya Bay's own real nominal free-normalization value (1.0).} & \pyfunc{load\_global\_normalization} \\
\bottomrule
\end{tabular}
\end{table}

\begin{notebox}
Not fetched/reproduced: the 6AD/7AD sub-periods (only 8AD); the
LSNL/IAV/background official \emph{uncertainty} parameters as a full
correlated covariance (their real nominal/central-value curves and rates
\emph{are} used, see \cref{tab:dayabay-tables-response}); and per-detector/per-hall
correlation structure for the background-rate systematics (a single
correlated scale per category is used instead, \cref{ch:inference-process}).
\end{notebox}


% =============================================================================
\chapter{Detector Implementation in \texttt{tpeanuts}}
\label{ch:detector-implementation}

This chapter develops, in full mathematical detail, every model
\texttt{tpeanuts.detector.dayabay} implements: the real multi-reactor
antineutrino flux (\cref{sec:flux-model}), the real order-1/M inverse-beta-decay
cross section (\cref{sec:ibd-cross-section}), the real detector energy
response (\cref{sec:iav-lsnl-theory}), efficiency and exposure
(\cref{sec:efficiency-exposure}), real backgrounds
(\cref{sec:backgrounds}), and the final composed event-rate assembly
(\cref{sec:event-rate-assembly}) with its free-normalization and nuisance
parameters (\cref{sec:nuisances}).

\section{Reactor antineutrino flux}
\label{sec:flux-model}

\modulehead{detector/dayabay/flux.py}{Real per-isotope Huber-Mueller
spectrum, with real non-equilibrium and spent-nuclear-fuel corrections, a
real per-reactor relative power weight, and geometric dilution.}

\begin{theorybox}
\subsection*{Per-isotope spectrum and fission-fraction sum}
The antineutrino spectrum shape at reactor core $r$ is a real
fission-fraction-weighted sum over the four fissile isotopes, with the real
non-equilibrium correction $C_i^{\rm neq}(E)$ folded into each isotope's
tabulated spectrum $S_i(E)$ \emph{before} the sum (published for
$^{235}$U/$^{239}$Pu/$^{241}$Pu only -- no correction exists for
$^{238}$U, which is left uncorrected):
\begin{equation}
  \keyresult{\Sigma(E) = \sum_{i} f_i\, S_i(E)\,\bigl(1+C_i^{\rm neq}(E)\bigr).}
  \label{eq:dayabay-spectrum-shape}
\end{equation}
The non-equilibrium correction accounts for long-lived fission-daughter
beta decays that have not yet reached secular equilibrium with their
production rate; it is largest ($\sim6\%$) just above the IBD threshold and
falls to zero by $\sim4.5$~MeV.

\subsection*{Absolute normalisation and real per-reactor corrections}
The absolute flux at reactor core $r$, before geometric dilution, combines
the real thermal power, the real fission-fraction-weighted mean energy per
fission $\bar e = \sum_i f_i \bar e_i$, a real per-reactor relative
power/output weight $w_r$, and a real per-reactor spent-nuclear-fuel (SNF)
correction $C_r^{\rm SNF}(E)$ (residual antineutrino emission from stored,
previously used fuel assemblies near the core):
\begin{equation}
  \keyresult{\Phi_{{\rm face},r}(E) = \frac{P_{\rm th}}{\bar e}\;
    w_r\;\bigl(1+C_r^{\rm SNF}(E)\bigr)\;\Sigma(E).}
  \label{eq:dayabay-face-flux}
\end{equation}
$w_r$ is derived (not itself a tabulated quantity) from the real weekly
antineutrino-rate history, \cref{sec:weekly-rate-table} below.

\subsection*{Geometric dilution and the multi-reactor sum}
The flux incident on detector $d$ sums the (independently propagated)
contribution of all 6 real cores, each diluted by the real inverse-square
baseline law and multiplied by that reactor's own real oscillation
survival probability:
\begin{equation}
  \keyresult{\Phi_d(E) = \sum_{r=1}^{6}
    \frac{\Phi_{{\rm face},r}(E)}{4\pi L_{r,d}^2}\;
    P_{\bar\nu_e\bar\nu_e}(E,L_{r,d}).}
  \label{eq:dayabay-flux-sum}
\end{equation}
$P_{\bar\nu_e\bar\nu_e}(E,L)$ is the standard vacuum 3-flavour survival
probability, computed once per (reactor, detector) pair at that pair's own
real baseline (\cref{ch:inference-process}); no matter effect is applied
(reactor antineutrinos at these energies and baselines cross negligible
Earth matter). \Cref{eq:dayabay-flux-sum} is evaluated independently for
every one of the 8 real detectors -- each sees its own real 6-baseline
combination.
\end{theorybox}

\subsection*{Real per-reactor relative weight}
\label{sec:weekly-rate-table}

\begin{theorybox}
The official weekly \code{neutrino\_rate} table (\cref{ch:data-release})
gives each real reactor core's own antineutrino emission rate, resolved
week by week across the full 2011--2020 data-taking period, distinguishing
the 6AD/7AD/8AD sub-periods via an \code{n\_det} column. Restricting to the
real 8AD-period weeks and averaging,
\begin{equation}
  \keyresult{w_r = \frac{\left\langle\text{neutrino\_rate}_r\right\rangle_{\rm 8AD}}
    {\frac{1}{6}\sum_{r'} \left\langle\text{neutrino\_rate}_{r'}\right\rangle_{\rm 8AD}},}
  \label{eq:dayabay-relative-weight}
\end{equation}
a 6-reactor average of exactly 1 by construction. This real per-reactor
weight replaces the previous implicit assumption that all 6 cores
contribute identically to \labelcref{eq:dayabay-face-flux} -- the real
measured weights differ from 1 only at the $\sim1$--$2\%$ level (real
Table~\ref{tab:relative-weights-example} in \cref{ch:results}), a small but
genuine, real per-reactor correction.
\end{theorybox}

\begin{implbox}
\begin{itemize}
  \item \pyfunc{reactor\_spectrum\_shape(E\_grid\_MeV)} implements
    \labelcref{eq:dayabay-spectrum-shape} via a cached, corrected copy of
    the raw Huber-Mueller curves (\pyfunc{\_corrected\_hm\_curves}, built
    once with the real non-equilibrium correction interpolated onto each
    isotope's own tabulated energy grid, 0 outside its real defined range).
  \item \pyfunc{reactor\_flux\_at\_face(E\_grid\_MeV, reactor)} implements
    \labelcref{eq:dayabay-face-flux}, with the real SNF correction
    interpolated onto \code{E\_grid\_MeV} (0 outside its own real range).
  \item \pyfunc{flux\_at\_detector(E\_grid\_MeV, baseline\_km, reactor)}
    divides by $4\pi L^2$ (converting \code{baseline\_km} to cm) -- one
    real (reactor, detector) term of \labelcref{eq:dayabay-flux-sum};
    \pyfunc{detector.dayabay.event\_rate.ibd\_event\_rate}
    (\cref{sec:event-rate-assembly}) performs the sum over the 6 real
    reactors explicitly.
  \item \pyfunc{detector.dayabay.parameters.REACTOR\_RELATIVE\_WEIGHT}
    implements \labelcref{eq:dayabay-relative-weight} once, at import time,
    from the real cached weekly-rate table.
\end{itemize}
\end{implbox}

\section{The real order-1/M inverse-beta-decay cross section}
\label{sec:ibd-cross-section}

\modulehead{detector/interaction/inverse\_beta\_decay.py}{Vogel \&
Beacom's threshold-respecting, order-1/M differential IBD cross section
\parencite{VogelBeacom1999}, built by numerical $\cos\theta$ integration
using Daya Bay's own real published coupling constants.}

\begin{theorybox}
\subsection*{Zeroth order (no-recoil) baseline}
At zeroth order in $1/M$ ($M$ the average nucleon mass), the positron
energy is fixed by $E_\nu$ alone,
\begin{equation}
  E_e^{(0)} = E_\nu - \Delta,\qquad \Delta = M_n - M_p,
\end{equation}
and the total cross section is
\begin{equation}
  \sigma^{(0)}(E_\nu) = \sigma_0\,(f^2+3g^2)\,E_e^{(0)} p_e^{(0)},\qquad
  p_e^{(0)} = \sqrt{\bigl(E_e^{(0)}\bigr)^2-m_e^2}.
  \label{eq:ibd-zeroth-order}
\end{equation}
This is Vogel \& Beacom's own Eq.~(10)-(11), and is the formula this
project used for every detector prior to this Daya Bay extension
(\texttt{sigma\_ibd}); it neglects recoil, weak magnetism, and the
resulting positron-angle dependence, and its reaction threshold,
$E_\nu > \bigl((M_n+m_e)^2-M_p^2\bigr)/(2M_p)\approx1.806$~MeV, is exact
only at this order.

\subsection*{Order-1/M differential cross section}
Vogel \& Beacom's own Eq.~(18)-(19) ``high-energy limit'' formula
explicitly neglects the threshold, an effect their own paper states is
``large'' below $\sim30$~MeV -- \emph{every} reactor antineutrino is below
10~MeV, so that formula is not used here. Instead, the full order-1/M
formula (their Eq.~13-15), valid down to threshold, is implemented
directly. At first order the positron energy depends on the
positron-antineutrino scattering angle $\theta$ (lab frame):
\begin{equation}
  \keyresult{E_e^{(1)} = E_e^{(0)}\left[1-\frac{E_\nu}{M}\bigl(1-v_e^{(0)}\cos\theta\bigr)\right]
    - \frac{y^2}{M},\qquad y^2=\frac{\Delta^2-m_e^2}{2},\qquad M=\frac{M_n+M_p}{2},}
  \label{eq:ibd-positron-energy-firstorder}
\end{equation}
with $v_e^{(0)}=p_e^{(0)}/E_e^{(0)}$, and $p_e^{(1)}=\sqrt{(E_e^{(1)})^2-m_e^2}$
(zero if the argument is negative, the exact real threshold cutoff at this
order). The differential cross section becomes
\begin{equation}
  \keyresult{\left(\frac{d\sigma}{d\cos\theta}\right)^{(1)} =
    \frac{\sigma_0}{2}\Bigl[(f^2+3g^2)+(f^2-g^2)\,v_e^{(1)}\cos\theta\Bigr]E_e^{(1)}p_e^{(1)}
    \;-\;\frac{\sigma_0}{2}\,\frac{\Gamma}{M}\,E_e^{(0)}p_e^{(0)},}
  \label{eq:ibd-firstorder-diffxs}
\end{equation}
with $\Gamma$ the weak-magnetism/recoil correction term,
\begin{align}
  \Gamma ={}& 2(f+f_2)g\Bigl[(2E_e^{(0)}+\Delta)(1-v_e^{(0)}\cos\theta)-\tfrac{m_e^2}{E_e^{(0)}}\Bigr] \nonumber\\
  &+ (f^2+g^2)\Bigl[\Delta(1+v_e^{(0)}\cos\theta)+\tfrac{m_e^2}{E_e^{(0)}}\Bigr] \nonumber\\
  &+ (f^2+3g^2)\Bigl[(E_e^{(0)}+\Delta)\bigl(1-\tfrac{\cos\theta}{v_e^{(0)}}\bigr)-\Delta\Bigr] \nonumber\\
  &+ (f^2-g^2)\Bigl[(E_e^{(0)}+\Delta)\bigl(1-\tfrac{\cos\theta}{v_e^{(0)}}\bigr)-\Delta\Bigr]v_e^{(0)}\cos\theta.
  \label{eq:ibd-gamma}
\end{align}

\subsection*{Removing the apparent threshold singularity}
\Cref{eq:ibd-gamma} contains terms $\propto 1/v_e^{(0)}$, apparently
divergent as $E_\nu\to$ threshold ($v_e^{(0)}\to0$); Vogel \& Beacom
themselves note this is exactly cancelled by the explicit $p_e^{(0)}$
prefactor multiplying $\Gamma$ in \labelcref{eq:ibd-firstorder-diffxs}.
Substituting $v_e^{(0)}=p_e^{(0)}/E_e^{(0)}$ throughout \emph{before}
multiplying by $p_e^{(0)}$ makes this cancellation manifest algebraically,
avoiding a numerically risky $0\times\infty$ evaluation at floating-point
precision:
\begin{align}
  p_e^{(0)}\Gamma ={}& 2(f+f_2)g\left[p_e^{(0)}(2E_e^{(0)}+\Delta) -
    \cos\theta\,\frac{\bigl(p_e^{(0)}\bigr)^2}{E_e^{(0)}}(2E_e^{(0)}+\Delta)
    - \frac{p_e^{(0)}m_e^2}{E_e^{(0)}}\right] \nonumber \\
  & + (f^2+g^2)\left[\Delta p_e^{(0)} + \cos\theta\,\Delta\,\frac{\bigl(p_e^{(0)}\bigr)^2}{E_e^{(0)}}
    + \frac{p_e^{(0)}m_e^2}{E_e^{(0)}}\right] \nonumber \\
  & + (f^2+3g^2)\Bigl[p_e^{(0)}E_e^{(0)}-\cos\theta\,E_e^{(0)}(E_e^{(0)}+\Delta)\Bigr] \nonumber \\
  & + (f^2-g^2)\,p_e^{(0)}\cos\theta\Bigl[p_e^{(0)}-\cos\theta(E_e^{(0)}+\Delta)\Bigr],
  \label{eq:ibd-gamma-clean}
\end{align}
free of any $1/v_e^{(0)}$ or $1/p_e^{(0)}$ term -- every division is by
$E_e^{(0)}$, which is bounded below by $m_e>0$ over the whole real IBD
kinematic range.

\subsection*{Bare normalisation $\sigma_0$}
The bare $\sigma_0$ (Vogel \& Beacom's Eq.~9) does not depend on
$f/g/f_2$; it is computed here directly rather than back-derived from the
zeroth-order coefficient (which was calibrated using Vogel \& Beacom's own,
slightly older $g=1.26$ value, not Daya Bay's real published
$g=1.2701$):
\begin{equation}
  \keyresult{\sigma_0 = \frac{G_F^2\cos^2\theta_C}{\pi}\bigl(1+\Delta_{\rm inner}^R\bigr)
    \times (\hbar c)^2,}
  \label{eq:ibd-sigma0-firstprinciples}
\end{equation}
with $\cos\theta_C=0.974$ the Cabibbo angle and $\Delta_{\rm inner}^R\approx
0.024$ the energy-independent inner radiative correction, both from
\cite{VogelBeacom1999}, and $(\hbar c)^2$ converting from natural units to
$\mathrm{cm^2/MeV^4}$. This reproduces Vogel \& Beacom's own quoted
$0.0952\times10^{-42}\ \mathrm{cm^2/MeV^2}$ coefficient to $\sim1\%$ when
recombined with their own $f=1$, $g=1.26$ -- the residual difference is
input-value vintage (rounding of $\Delta_{\rm inner}^R$/$\cos\theta_C$),
not a formula error.

\subsection*{Total cross section and the prompt-energy differential}
The total cross section is obtained by numerical integration over
$\cos\theta$,
\begin{equation}
  \sigma_{\rm tot}(E_\nu) = \int_{-1}^{1} d\cos\theta \left(\frac{d\sigma}{d\cos\theta}\right)^{(1)},
\end{equation}
the same method Vogel \& Beacom themselves used for their own published
curves. To build the differential cross section in the observable prompt
energy $d\sigma/dT$, each $\cos\theta$-quadrature node's cross-section
contribution is redistributed onto the prompt-energy grid at its own
order-1/M prompt energy,
\begin{equation}
  E_{\rm prompt} = E_e^{(1)} + m_e,
\end{equation}
via linear-interpolation mass conservation (identical technique to
\labelcref{eq:iav-warp-weights} below), then converted from a
per-quadrature-node mass to a density by dividing by the prompt-energy grid
spacing.
\end{theorybox}

\begin{implbox}
\begin{itemize}
  \item \pyfunc{sigma\_ibd\_diff\_cos\_precise(E\_nu\_MeV, cos\_theta, f=1.0,
    g=1.2701, f2=3.706)} implements
    \labelcref{eq:ibd-positron-energy-firstorder,eq:ibd-firstorder-diffxs,eq:ibd-gamma-clean}
    directly, defaulting to Daya Bay's own real coupling constants
    (\code{ibd\_constants.csv}, \cref{ch:data-release}); returns
    $(d\sigma/d\cos\theta,\,E_e^{(1)})$.
  \item \pyfunc{ibd\_cross\_section\_grid\_precise(E\_nu\_grid\_MeV,
    T\_grid\_MeV, n\_cos=41)} performs the $\cos\theta$-quadrature-and-scatter
    construction above on a uniform 41-point trapezoidal grid, returning
    $d\sigma/dT$ on the outer-product grid
    \code{(E\_nu\_grid\_MeV, T\_grid\_MeV)}, shape \code{(n\_E,n\_T)},
    $\mathrm{cm^2/MeV}$.
  \item \pyfunc{\_sigma0\_bare\_cm2\_per\_mev4()} implements
    \labelcref{eq:ibd-sigma0-firstprinciples}, cached at module import.
  \item \pyfunc{sigma\_ibd(E\_nu\_MeV)}/\pyfunc{ibd\_cross\_section\_grid(...)}:
    the original zeroth-order formula, \labelcref{eq:ibd-zeroth-order},
    retained unchanged (used by any detector that has not opted into the
    precise cross section).
  \item \pyfunc{detector.common.response.scatter\_add\_linear(grid,
    positions, values)}: the shared mass-conserving linear-interpolation
    scatter utility used both here (redistributing cross-section mass
    across $\cos\theta$ nodes) and by the LSNL warp matrix
    (\cref{sec:iav-lsnl-theory}).
\end{itemize}
\end{implbox}

\begin{notebox}
The real, measured effect of this upgrade: the order-1/M cross section is a
few percent \emph{below} the zeroth-order one, growing with $E_\nu$
(consistent with Vogel \& Beacom's own stated $\mathcal O(E_\nu/M)$
scaling, with an $\mathcal O(-7)$ numerical coefficient) -- validated
independently in this project's own test suite
(\code{detector/dayabay/test/test2\_response\_and\_fit.py}).
\end{notebox}

\section{Detector energy response: IAV, LSNL, and Gaussian resolution}
\label{sec:iav-lsnl-theory}

\modulehead{detector/dayabay/response.py}{The full real response chain,
applied in the official Daya Bay Collaboration analysis pipeline's own
order: IAV redistribution, then LSNL nonlinearity, then Gaussian
photostatistics resolution.}

\begin{theorybox}
\subsection*{Real response order}
The official Daya Bay dagflow/GNA analysis pipeline's own node names
(confirmed directly from the Collaboration's own CHEP~2026 conference
presentation on that framework) read, cumulatively,
\begin{equation}
  \text{raw} \;\longrightarrow\; \text{stages.iav} \;\longrightarrow\;
  \text{stages.evis (IAV+LSNL)} \;\longrightarrow\; \text{stages.erec (+resolution)},
\end{equation}
i.e.\ \keyresult{IAV is applied to the true deposited energy \emph{before}
LSNL}, which is then followed by the Gaussian resolution smearing. This
also matches the physical picture: energy lost into the inert acrylic
vessel material near its wall (the IAV effect) never produces
scintillation light, so it cannot be subject to a light-yield nonlinearity
curve (LSNL) -- only the light-producing, IAV-redistributed energy is.

\subsection*{IAV redistribution}
The real IAV correction is a $240\times240$ energy-redistribution matrix
$R^{\rm IAV}$ on a $0.05$~MeV, $0$--$12$~MeV grid (matching this project's
own \code{T\_GRID\_MEV} exactly), normalised so that each column sums to
1 -- a genuine discrete probability-mass (not density) transfer matrix,
\begin{equation}
  R^{\rm IAV}_{ij} = P(\text{reconstructed in bin } i \mid \text{true energy in bin } j),
  \qquad \sum_i R^{\rm IAV}_{ij} = 1.
\end{equation}

\subsection*{LSNL nonlinearity}
The real LSNL correction is a multiplicative energy-scale warp,
$f(E)=E_{\rm reconstructed}/E_{\rm true\ deposited}$, mapping true energy
$T$ to $T_{\rm nl}=T\cdot f(T)$. Represented as a mass-conserving matrix on
the same grid, column $j$ (true energy $T_j$) is split by linear
interpolation between the two grid points bracketing $T_{\rm nl}(T_j)$:
\begin{equation}
  \keyresult{R^{\rm LSNL}_{ij} = \begin{cases}
    1-w_j, & i = \lfloor T_{\rm nl}(T_j)\rfloor_{\rm grid}, \\
    w_j, & i = \lceil T_{\rm nl}(T_j)\rceil_{\rm grid}, \\
    0, & \text{otherwise,}
  \end{cases}}
  \qquad w_j = \frac{T_{\rm nl}(T_j)-T_{\lfloor\cdot\rfloor}}{T_{\lceil\cdot\rceil}-T_{\lfloor\cdot\rfloor}}
  \label{eq:iav-warp-weights}
\end{equation}
(clamped to $[0,1]$, with all mass dumped at the nearest edge for a
position outside the grid range) -- the identical construction used for
the IBD cross section's $\cos\theta\to T$ redistribution
(\cref{sec:ibd-cross-section}). By construction $\sum_i R^{\rm LSNL}_{ij}=1$
for every column $j$.

\subsection*{Official LSNL pull-curve model}
The real LSNL curve used is not fixed to its nominal shape alone: the
official systematic model is a linear combination of the nominal curve
$f_0$ and 4 real published pull curves $f_k$,
\begin{equation}
  \keyresult{f(E) = f_0(E) + \sum_{k=1}^{4} a_k\bigl(f_k(E)-f_0(E)\bigr),}
  \label{eq:lsnl-pull-model}
\end{equation}
with each $a_k$ a real free nuisance carrying an official independent,
zero-mean, unit-sigma Gaussian prior (\code{parameters/detector\_lsnl.yaml}).
At $a_k=0$ for all $k$, \labelcref{eq:lsnl-pull-model} reduces to the
nominal curve; \cref{ch:inference-process} documents how $a_k$ enters the
fit as a real, Gaussian-constrained nuisance parameter.

\subsection*{Gaussian photostatistics resolution}
The real 3-term resolution formula,
\begin{equation}
  \keyresult{\frac{\sigma(E)}{E} = \sqrt{a^2+\frac{b^2}{E}+\frac{c^2}{E^2}}},
  \qquad a=0.016,\ b=0.081\ \sqrt{\rm MeV},\ c=0.026\ {\rm MeV},
  \label{eq:eres-formula}
\end{equation}
with terms interpretable as spatial/temporal non-uniformity ($a$), photon
(photoelectron) counting statistics ($b$), and PMT dark noise/electronics
($c$), builds a genuine probability \emph{density} in $T'$ (not a discrete
mass matrix, unlike IAV/LSNL): $R^{\rm Gauss}(T'|T) =
\mathcal N(T';T,\sigma(T))$.

\subsection*{Composition}
All three steps are real, independent, official detector-response
ingredients (the data release ships them as separate real files precisely
because they are physically distinct effects); they are composed as
sequential linear operators. Because the folding machinery
(\cref{sec:event-rate-assembly}) integrates $R(T'|T)$ against $T$ via
trapezoidal quadrature, every step is first built as a discrete,
mass-conserving transfer matrix on the shared 0.05~MeV grid, and only the
final \emph{composed} matrix is converted from probability mass to the
density the folding integral expects (dividing by the grid spacing $\delta
T$):
\begin{equation}
  \keyresult{R(T'|T) = \frac{1}{\delta T}\,
    R^{\rm Gauss}_{\rm mass}\;R^{\rm LSNL}\;R^{\rm IAV},}
  \label{eq:response-composition}
\end{equation}
where $R^{\rm Gauss}_{\rm mass} = \delta T\cdot R^{\rm Gauss}$ (density
converted to an approximate mass for the composition, valid since $\delta
T=0.05$~MeV is fine compared to $\sigma(E)$ over the whole relevant
energy range). Reading \labelcref{eq:response-composition} as a matrix
product applied to a column vector $v$ representing a true-energy
spectrum, $R(T'|T)\,v = R^{\rm Gauss}_{\rm mass}\bigl(R^{\rm LSNL}(R^{\rm
IAV}v)\bigr)/\delta T$ -- $v$ is transformed by $R^{\rm IAV}$ first, then
$R^{\rm LSNL}$, then $R^{\rm Gauss}_{\rm mass}$, exactly the real physical
order stated above.
\end{theorybox}

\begin{implbox}
\begin{itemize}
  \item \pyfunc{sigma\_MeV(T\_grid\_MeV, a, b, c)} implements
    \labelcref{eq:eres-formula}.
  \item \pyfunc{\_iav\_mass\_matrix(T\_grid\_MeV)} embeds the real
    $240\times240$ \pyfunc{IAV\_MATRIX} into the project's own (possibly
    larger) grid, with an identity fallback for any extra grid point beyond
    the real IAV table's own range (physically negligible: the real IBD
    prompt-energy spectrum is essentially zero there).
  \item \pyfunc{\_lsnl\_curves\_on\_grid(T\_grid\_MeV)} interpolates the
    real nominal/pull curves onto the shared grid once
    (\code{@torch.no\_grad()}, fixed real data).
  \item \pyfunc{\_lsnl\_warp\_matrix(T\_grid\_MeV, lsnl\_pulls=None)}
    implements \labelcref{eq:lsnl-pull-model,eq:iav-warp-weights};
    differentiable w.r.t.\ \code{lsnl\_pulls} (only the linear combination
    and the resulting warp weights carry gradient -- the underlying curve
    interpolation is fixed real data).
  \item \pyfunc{response\_matrix(T\_grid\_MeV, Tprime\_grid\_MeV, a, b, c,
    lsnl\_pulls=None)} implements \labelcref{eq:response-composition}.
\end{itemize}
\end{implbox}

\begin{notebox}
Every component above uses the real official nominal/central-value
curves/matrix; the LSNL pull-curve nuisances $a_k$ can be fit
(\cref{ch:inference-process}), but the IAV matrix's own official
off-diagonal-scale systematic (\code{detector\_iav\_offdiag\_scale.yaml},
a single real nuisance with a $4\%$ prior) is not implemented. The three
response steps are composed assuming they act as independent, sequential
linear operators -- a standard simplification when a full joint response
covariance is not being reproduced.
\end{notebox}

\section{Efficiency and exposure}
\label{sec:efficiency-exposure}

\begin{theorybox}
The real IBD-selection efficiency $\varepsilon\approx0.8025$
(Gd-capture fraction $\times$ analysis-cut efficiencies, consistent with
the real published individual cut efficiencies -- Gd-capture fraction
$\approx85.5\%$, prompt-energy cut $\approx99.81\%$, capture-time cut
$\approx98.70\%$, multiplicity cut $\approx97.5\%$, flasher cut
$\approx99.98\%$, whose product is $\approx81.6\%$, close to the quoted
$80.25\%$) is applied as a flat multiplicative factor, on top of and
separate from the real per-detector, per-day effective livetime
(\code{eff\_livetime} $=$ \code{livetime}$\times$\code{eff}, the latter
itself the real combined muon-veto and multiplicity-cut daily efficiency),
summed over the real 8AD-period days to give each detector's total real
exposure in seconds.
\end{theorybox}

\begin{implbox}
\pyfunc{DETECTOR\_EFFICIENCY} (a flat scalar) multiplies the folded
reconstructed spectrum after response and before binning
(\cref{sec:event-rate-assembly}); \pyfunc{load\_exposure()} sums
\code{eff\_livetime} over real 8AD-period days per detector.
\end{implbox}

\section{Real backgrounds}
\label{sec:backgrounds}

\begin{theorybox}
Five real background categories contribute to the observed spectrum,
each with its own real published per-detector rate (events/day) and real
normalised shape histogram (integrating to 1 on the same 0.05~MeV binning
as the signal):
\begin{itemize}
  \item \textbf{Accidentals} -- uncorrelated prompt-like/delayed-like
    triggers randomly coincident in time; the real best-measured category
    (fractional rate uncertainty $\sim0.1\%$).
  \item \textbf{$^9$Li/$^8$He} -- cosmogenic isotopes produced by
    through-going muons, beta-decaying with a delayed-neutron-emission
    branch that mimics the IBD signature.
  \item \textbf{Fast neutrons} -- energetic neutrons from muon spallation
    outside the water shield, which can deposit a prompt-like recoil energy
    followed by a real neutron capture.
  \item \textbf{AmC} -- correlated background from the AmC neutron
    calibration sources themselves.
  \item \textbf{$^{13}$C$(\alpha,n)^{16}$O} -- alpha decays of natural
    radioactivity in the scintillator inducing a real $(\alpha,n)$ reaction
    on carbon.
\end{itemize}
The predicted background count per real analysis bin is
\begin{equation}
  N_k^{\rm bkg} = \sum_{\rm category} ({\rm rate}_{\rm category}\times{\rm exposure\_days})
    \times{\rm shape}_{{\rm category},k},
\end{equation}
with the fine-binned (0.05~MeV) product rebinned onto the real
non-uniform analysis edges by a plain per-bin sum (every real analysis-bin
edge lands exactly on a fine-bin boundary, so no fine bin straddles two
coarse bins).
\end{theorybox}

\begin{implbox}
\pyfunc{real\_background\_counts(detector, exposure\_seconds,
bin\_edges\_MeV, category\_scale=None)} implements the sum above;
\pyfunc{rebin\_discrete\_counts} performs the fine-to-coarse rebinning.
An optional \code{category\_scale} dict multiplies each category's
contribution by a real per-category nuisance (\cref{sec:nuisances}).
\end{implbox}

\section{Composed event-rate assembly}
\label{sec:event-rate-assembly}

\modulehead{detector/dayabay/event\_rate.py}{The one place this package's
real 6-reactor geometry, real flux, real order-1/M cross section, real
response, and real background come together, per detector.}

\begin{theorybox}
For one real detector $d$, the predicted signal (pre-response) spectrum
folds the multi-reactor flux sum \labelcref{eq:dayabay-flux-sum} with the
order-1/M differential cross section \labelcref{eq:ibd-firstorder-diffxs},
\begin{equation}
  \frac{dR}{dT}(T) = N_p^{(d)}\int dE_\nu\;\Phi_d(E_\nu)\,\frac{d\sigma}{dT}(E_\nu,T),
\end{equation}
with $N_p^{(d)}$ the real target free-proton count for detector $d$; the
response matrix \labelcref{eq:response-composition} is then applied,
\begin{equation}
  \frac{dR}{dT'}(T') = \int dT\; R(T'|T)\,\frac{dR}{dT}(T),
\end{equation}
followed by the flat efficiency factor and integration into the real
non-uniform analysis bins,
\begin{equation}
  \keyresult{N_k(\theta) = \nu_{\rm norm}\cdot{\rm exposure}\cdot\varepsilon\int_{{\rm bin}\,k}
    dT'\;\frac{dR}{dT'}(T';\theta) \;+\; N_k^{\rm bkg},}
  \label{eq:dayabay-full-prediction}
\end{equation}
the full, real, per-bin, per-detector prediction, with $\nu_{\rm norm}$
the free global-normalization nuisance (\cref{sec:nuisances}) multiplying
the signal only (the real background is separately measured and not
rescaled by it).
\end{theorybox}

\begin{implbox}
\begin{itemize}
  \item \pyfunc{ibd\_event\_rate(detector, p\_ee\_per\_reactor, ...,
    signal\_scale=1.0, background\_category\_scale=None, lsnl\_pulls=None)}
    implements \labelcref{eq:dayabay-full-prediction} end to end, calling
    \pyfunc{detector.common.event\_rate.predicted\_counts} with the real
    cross section, real response, real target, real exposure, and real
    background plugged in.
  \item \code{p\_ee\_per\_reactor}: \code{\{reactor: P\_ee(E)\}}, one
    tensor per real reactor, already oscillated at that reactor's own
    real baseline -- supplied by
    \pyclass{detector.dayabay.inference\_model.DayaBayDetectorModel}
    (\cref{ch:inference-process}), which itself calls
    \texttt{tpeanuts.medium.vacuum.oscillation\_model
    .VacuumOscillationModel.predict\_pee} once per reactor.
\end{itemize}
\end{implbox}

\section{Free normalization and nuisance parameters}
\label{sec:nuisances}

\begin{theorybox}
\subsection*{Global normalization}
Daya Bay's own official combined analysis leaves a free overall
normalization nuisance in its fit (\code{parameters/detector\_normalization.yaml},
real nominal value 1.0, labelled ``Global IBD rate normalization''): even
Daya Bay's own simulation chain is not asserted to self-normalize to the
observed rate without fitting it. \Cref{eq:dayabay-full-prediction}'s
$\nu_{\rm norm}$ is this same real, official parameter, not an ad hoc
rescaling invented for this project -- see \cref{ch:results} for why it is
needed (a real, documented $\sim19\%$ absolute-normalization gap, audited
and explained in \cref{ch:results}, not attributable to any single unit or
loading error).

\subsection*{Real background-rate nuisances}
Five further real nuisances, one Gaussian-constrained multiplicative scale
$c_k$ per background category (nominal 1), with real prior width
\begin{equation}
  \sigma_k = \frac{\sqrt{\sum_d ({\rm uncertainty}_{k,d})^2}}{\sum_d {\rm rate}_{k,d}},
  \label{eq:background-sigma}
\end{equation}
a rate-weighted quadrature combination of the real per-detector
(rate, uncertainty) pairs into one correlated scale per category -- a
simplification of the official per-detector/per-hall correlation structure
(accidentals genuinely uncorrelated between detectors; AmC correlated
across all detectors; $^9$Li/$^8$He and fast neutrons correlated within a
hall), not a per-detector free parameter.

\subsection*{Real LSNL pull-curve nuisances}
The 4 real $a_k$ of \labelcref{eq:lsnl-pull-model}, official prior
independent zero-mean unit-sigma Gaussian.
\end{theorybox}

\begin{implbox}
\begin{itemize}
  \item \pyclass{DayaBayDetectorModel}(\code{normalization\_free},
    \code{background\_free}, \code{lsnl\_free}): each flag appends the
    corresponding real nuisance name(s) to \code{.free} and slices them
    from \code{theta} in \code{.predict} (\cref{ch:inference-process}).
  \item \pyfunc{detector.dayabay.parameters.BACKGROUND\_CATEGORY\_SIGMA}
    implements \labelcref{eq:background-sigma} once, at import time.
  \item \pyfunc{tpeanuts.inference.likelihood.gaussian\_prior\_penalty(values,
    mean, sigma)}: the generic, reusable $-2\ln L$ Gaussian-prior penalty
    added to the fit's likelihood for any of these nuisances
    (\cref{ch:inference-process}).
\end{itemize}
\end{implbox}

 
% =============================================================================
\chapter{The Inference Process}
\label{ch:inference-process}

This chapter documents the statistical inference machinery: the 3-flavour
vacuum oscillation probability every model below shares
(\cref{sec:oscillation-probability}), the three composed model variants
(\cref{sec:model-variants}), the likelihood statistics
(\cref{sec:likelihoods}), the gradient-based fit procedure and its
uncertainty estimate (\cref{sec:fit-procedure}), and the 2-D likelihood
scan (\cref{sec:likelihood-scan}).

\section{Real 3-flavour vacuum oscillation probability}
\label{sec:oscillation-probability}

\modulehead{medium/vacuum/oscillation\_model.py}{Shared, detector-agnostic
3-flavour vacuum oscillation model (\pyclass{VacuumOscillationModel}),
reused unchanged from the rest of \texttt{tpeanuts}.}

\begin{theorybox}
Reactor antineutrinos at Daya Bay's real baselines ($\lesssim2$~km) cross
negligible Earth matter, so the exact coherent 3-flavour \emph{vacuum}
survival probability is used,
\begin{equation}
  P_{\bar\nu_e\bar\nu_e}(E,L;\theta) = \left|\sum_i U_{ei}^{*}\,
    e^{-i\Delta m^2_{i1}L/2E}\,U_{ei}\right|^2,
  \label{eq:vacuum-survival}
\end{equation}
with $U$ the PMNS matrix \parencite{Pontecorvo1957,MakiNakagawaSakata1962}
and antineutrino sign convention ($U\to U^*$ relative to the neutrino
case). Free parameters throughout this document: $\theta_{13}$ and $\Delta
m^2_{3l}\equiv\Delta m^2_{31}$ (normal ordering); fixed at Daya Bay's own
real external solar input, $\theta_{12}=\tfrac12\arcsin\sqrt{\sin^2
2\theta_{12}}$, $\Delta m^2_{21}$ (\code{survival\_probability\_solar.yaml},
\cref{ch:data-release}). $\theta_{23}$/$\delta_{CP}$ do not affect
\labelcref{eq:vacuum-survival} (a direct consequence of the standard
3-flavour PMNS parametrisation) and are held at an arbitrary NuFIT preset
value \parencite{EstebanNuFit2024}.
\end{theorybox}

\begin{implbox}
\pyclass{VacuumOscillationModel.predict\_pee(theta, L\_km, E\_grid\_MeV)}
implements \labelcref{eq:vacuum-survival} exactly (full 3-flavour, not a
2-flavour approximation), called once per (reactor, detector) pair by
\pyclass{DayaBayDetectorModel.predict} (\cref{sec:model-variants}).
\end{implbox}

\section{The three model variants}
\label{sec:model-variants}

\modulehead{detector/dayabay/inference\_model.py}{\pyclass{DayaBayDetectorModel},
\pyclass{JointDayaBayModel}, \pyclass{NearFarRatioDayaBayModel},
\pyclass{DayaBayExperimentalModel}: composition layer wrapping
\pyclass{VacuumOscillationModel} around \texttt{event\_rate.ibd\_event\_rate}
(\cref{ch:detector-implementation}).}

\subsection*{Per-detector and joint absolute-count models}

\begin{theorybox}
\pyclass{DayaBayDetectorModel} evaluates \labelcref{eq:vacuum-survival}
once per real reactor (6 calls, each at that reactor's own real baseline),
then calls \labelcref{eq:dayabay-full-prediction} to fold the real
6-reactor flux sum through the real cross section, response, and
background. Its free-parameter vector is assembled, in order, from
\begin{equation}
  \theta = \bigl(\underbrace{\theta_{13},\,\Delta m^2_{3l}}_{\rm oscillation}\bigr)
  \oplus \bigl[\nu_{\rm norm}\bigr]
  \oplus \bigl[c_1,\dots,c_5\bigr]
  \oplus \bigl[a_1,\dots,a_4\bigr],
  \label{eq:theta-assembly}
\end{equation}
with each bracketed group present only if the corresponding
\code{normalization\_free}/\code{background\_free}/\code{lsnl\_free} flag
is set. \pyclass{JointDayaBayModel} concatenates several
\pyclass{DayaBayDetectorModel} predictions (typically all 8 real
detectors) sharing one $\theta$, for a full combined fit -- the
\textbf{full-spectral-shape} strategy.
\end{theorybox}

\subsection*{Rate-only model}

\begin{theorybox}
A rate-only observable integrates \labelcref{eq:dayabay-full-prediction}
over the full real analysis-bin range, discarding the intra-detector
spectral shape but retaining the real detector-to-detector rate deficit:
\begin{equation}
  \keyresult{N_d^{\rm rate}(\theta) = \sum_k N_k^{(d)}(\theta),}
  \label{eq:rate-only-observable}
\end{equation}
one real number per detector. Because a rate-only observable has
essentially no independent $\Delta m^2_{31}$ leverage in this two-baseline-tier
geometry, $\Delta m^2_{31}$ is fixed at Daya Bay's own published value and
only $\theta_{13}$ (plus $\nu_{\rm norm}$) is fit.
\end{theorybox}

\subsection*{Near/far ratio model}

\begin{theorybox}
Mirroring Daya Bay's own original 2012 discovery strategy
\parencite{An2012DayaBayDiscovery}, the far-hall/near-hall per-bin count
ratio is predicted instead of absolute counts:
\begin{equation}
  \keyresult{R_k(\theta) = \frac{N_k^{\rm far}(\theta)}{N_k^{\rm near}(\theta)},
    \qquad N_k^{\rm near,far}(\theta) = \!\!\sum_{d\,\in\,{\rm NEAR,FAR}}\!\! N_k^{(d)}(\theta),}
  \label{eq:nearfar-ratio}
\end{equation}
with $\mathrm{NEAR}=\{\mathrm{AD11,AD12,AD21,AD22}\}$ (EH1+EH2) and
$\mathrm{FAR}=\{\mathrm{AD31,\dots,AD34}\}$ (EH3,
\cref{sec:detector-geometry}). This cancels, by construction, every
systematic common to all 8 real detectors: with \code{normalization\_free
=False}, $\nu_{\rm norm}$ is not even a free parameter of this model (it
would cancel in \labelcref{eq:nearfar-ratio} to the extent the real
background is negligible relative to signal, and only approximately once
it is not); any \emph{residual} mismatch in the simplified IAV/LSNL/cross-section
response also cancels approximately, since it multiplies the near and far
prediction identically.
\end{theorybox}

\subsection*{Experimental reparametrisation}

\begin{theorybox}
Daya Bay's own combined analyses report $\sin^2(2\theta_{13})$ and the
effective splitting
\begin{equation}
  \Delta m^2_{ee} = \cos^2\theta_{12}\,\Delta m^2_{31} + \sin^2\theta_{12}\,\Delta m^2_{32}
  \label{eq:dm2ee-definition}
\end{equation}
\parencite{NunokawaParkeZukanovich2005,An2023DayaBayPrecision}, not the bare
PMNS $(\theta_{13},\Delta m^2_{3l})$ this project's models natively expose.
With $\theta_{12}$/$\Delta m^2_{21}$ fixed, substituting $\Delta
m^2_{32}=\Delta m^2_{31}-\Delta m^2_{21}$ into
\labelcref{eq:dm2ee-definition} and solving gives an exact, fixed linear
offset,
\begin{equation}
  \keyresult{\Delta m^2_{31} = \Delta m^2_{ee} + \sin^2\theta_{12}\,\Delta m^2_{21},}
  \label{eq:dm31-from-dm2ee}
\end{equation}
so \pyclass{DayaBayExperimentalModel} is a pure, exactly invertible change
of variables (verified in this project's own test suite), not a different
fit -- it exists purely so a fit's free-parameter vector and any resulting
contour are directly comparable to Daya Bay's own published numbers/plots.
\end{theorybox}

\begin{implbox}
\begin{itemize}
  \item \pyclass{DayaBayDetectorModel(oscillation\_model, detector,
    normalization\_free=True, background\_free=False, lsnl\_free=False)}:
    \code{.free} implements \labelcref{eq:theta-assembly};
    \code{.predict(theta)} slices $\theta$ in that same order and calls
    \pyfunc{ibd\_event\_rate}.
  \item \pyclass{JointDayaBayModel(models)}: \code{.predict} concatenates.
  \item \pyclass{NearFarRatioDayaBayModel(near\_models, far\_models)}
    implements \labelcref{eq:nearfar-ratio};
    \code{.real\_observed\_ratio()} (a \code{@staticmethod}) returns the
    real observed ratio and its propagated uncertainty
    (\cref{sec:likelihoods}).
  \item \pyclass{DayaBayExperimentalModel(model, theta12, dm21)} implements
    \labelcref{eq:dm31-from-dm2ee}, converting \code{theta}'s first two
    entries before delegating to the wrapped model.
\end{itemize}
\end{implbox}

\section{Likelihood statistics}
\label{sec:likelihoods}

\modulehead{inference/likelihood.py}{\pyfunc{poisson\_nll},
\pyfunc{chi2\_asymmetric}, \pyfunc{gaussian\_prior\_penalty} -- every
$-2\ln L$ statistic used in this document.}

\begin{theorybox}
\subsection*{Poisson (Baker-Cousins) likelihood}
Absolute per-bin counts (full-shape and rate-only fits) are compared via
the Poisson likelihood-ratio statistic \parencite{BakerCousins1984},
\begin{equation}
  \keyresult{-2\ln L_{\rm Poisson} = 2\sum_k\Bigl[N_k-n_k^{\rm obs}+n_k^{\rm obs}
    \ln\frac{n_k^{\rm obs}}{N_k}\Bigr],}
  \label{eq:poisson-nll}
\end{equation}
with $N_k$/$n_k^{\rm obs}$ the predicted/observed counts. The extra
$-n_k^{\rm obs}+n_k^{\rm obs}\ln n_k^{\rm obs}$ terms subtract the
saturated model's own likelihood, so \labelcref{eq:poisson-nll} is exactly
0 when $N_k=n_k^{\rm obs}$ bin by bin and behaves like an ordinary
chi-square asymptotically (Wilks' theorem).

\subsection*{Gaussian (chi2\_asymmetric) likelihood}
The near/far ratio (a ratio of counts, not itself Poisson-distributed) is
compared via a Gaussian statistic with one-sided uncertainty,
\begin{equation}
  -2\ln L_{\chi^2} = \sum_k \left(\frac{R_k-R_k^{\rm obs}}{\sigma_k}\right)^2,
  \label{eq:chi2-asymmetric}
\end{equation}
with the real observed ratio's uncertainty propagated from independent
Poisson numerator/denominator counts,
\begin{equation}
  \keyresult{\sigma_k = R_k^{\rm obs}\sqrt{\frac{1}{N_k^{\rm far,obs}}+\frac{1}{N_k^{\rm near,obs}}}.}
  \label{eq:ratio-sigma-propagation}
\end{equation}

\subsection*{Gaussian prior penalty}
Real, externally measured nuisance-parameter constraints (background-rate
scales, LSNL pulls, \cref{sec:nuisances}) are added as an independent
Gaussian penalty on the free-parameter vector itself, not the prediction:
\begin{equation}
  \keyresult{-2\ln L_{\rm penalty} = \sum_j \left(\frac{\theta_j-\mu_j}{\sigma_j}\right)^2,}
  \label{eq:gaussian-prior-penalty}
\end{equation}
added to \labelcref{eq:poisson-nll} (or \labelcref{eq:chi2-asymmetric}) to
give the total statistic minimized.
\end{theorybox}

\begin{implbox}
\begin{itemize}
  \item \pyfunc{poisson\_nll(prediction, value)} implements
    \labelcref{eq:poisson-nll}, built as
    \code{xlogy(value,value)-xlogy(value,prediction)} rather than the
    algebraically equal single-\code{xlogy} form, to keep the backward
    pass finite at an empty bin (\code{value==0}).
  \item \pyfunc{chi2\_asymmetric(prediction, value, sigma\_minus, sigma\_plus)}
    implements \labelcref{eq:chi2-asymmetric}.
  \item \pyfunc{gaussian\_prior\_penalty(values, prior\_mean, prior\_sigma)}
    implements \labelcref{eq:gaussian-prior-penalty}; passed as
    \pyfunc{fit\_lbfgs}'s \code{penalty\_fn} argument
    (\cref{sec:fit-procedure}).
  \item \pyfunc{NearFarRatioDayaBayModel.real\_observed\_ratio} implements
    \labelcref{eq:ratio-sigma-propagation}.
\end{itemize}
\end{implbox}

\section{Gradient-based fit procedure}
\label{sec:fit-procedure}

\modulehead{inference/fit.py}{\pyfunc{fit\_lbfgs}/\pyfunc{minimize\_lbfgs}:
LBFGS point estimate plus Laplace/Fisher covariance, for any
\pyclass{DifferentiableModel}.}

\begin{theorybox}
\subsection*{Parameter rescaling}
LBFGS's identity-initialised inverse-Hessian approximation assumes roughly
comparable parameter scales; $\theta_{13}\sim\mathcal O(1)$~rad and
$\Delta m^2_{31}\sim\mathcal O(10^{-3})$~eV$^2$ differ by 3-4 orders of
magnitude, which starves the line search of a usable step size. A
dimensionless $u=\theta/{\rm scale}$ is fit instead, with ${\rm scale}$ set
from the loss gradient at $\theta_0$,
\begin{equation}
  {\rm scale} = \bigl|\nabla_\theta(-2\ln L)\bigr|_{\theta_0}^{-1},
\end{equation}
undone only in the returned $\hat\theta$ -- every model always sees
physical-unit $\theta$.

\subsection*{Optimisation and convergence}
$\hat\theta=\arg\min_\theta\bigl[-2\ln L(\theta) + {\rm penalty}(\theta)\bigr]$
is found by \code{torch.optim.LBFGS} using the exact autograd gradient of
the total statistic through every real physics step of
\cref{ch:detector-implementation} (no finite differences).

\subsection*{Laplace/Fisher covariance}
The marginal uncertainty on each fitted parameter is the Laplace
approximation: with $\mathcal I=\tfrac12\nabla^2_\theta(-2\ln L)$ the
Fisher information at the optimum (the autograd Hessian of the total,
penalised statistic),
\begin{equation}
  \keyresult{\Sigma = \mathcal I^{-1}, \qquad \sigma_j = \sqrt{\Sigma_{jj}}.}
  \label{eq:laplace-covariance}
\end{equation}
A direction the data barely constrains gives a Hessian eigenvalue near
zero (occasionally slightly negative from numerical round-off); small or
negative eigenvalues are floored at a tiny fraction of the largest one
before inversion, so an unconstrained direction gets a large but finite
$\sigma_j$ instead of \code{nan}.
\end{theorybox}

\begin{implbox}
\begin{itemize}
  \item \pyfunc{fit\_lbfgs(model, theta0, value, sigma\_minus=None,
    sigma\_plus=None, likelihood=..., penalty\_fn=None, max\_iter=200)}
    returns a \pyclass{FitResult} (\code{theta\_hat}, \code{param\_names},
    \code{chi2\_history}, \code{covariance}, and a \code{sigma} property).
  \item \code{penalty\_fn}: an optional \code{Callable[[theta],\ Tensor]}
    added to the loss inside the LBFGS closure \emph{and} to the Hessian
    computation for \labelcref{eq:laplace-covariance} -- a tight prior
    correctly shrinks its own parameter's reported uncertainty. \code{None}
    (default) reproduces the previous, unconstrained-fit behaviour exactly.
\end{itemize}
\end{implbox}

\section{2-D likelihood scan}
\label{sec:likelihood-scan}

\modulehead{inference/scan.py}{\pyfunc{loglik\_grid}: evaluates $-2\ln L$
on a $(\theta_x,\theta_y)$ grid, for visual/contour diagnostics.}

\begin{theorybox}
On a chosen 2-D grid, every entry of \code{model.free} other than the two
scanned parameters is held fixed at a baseline $\theta$ (a cheap slice
through the likelihood surface) or re-minimised at each grid point (the
statistically correct construction when other parameters are also
uncertain, at the cost of one LBFGS run per grid point). With exactly two
free parameters in \code{model.free} (the rate-only and near/far models,
once their own nuisances are fixed/absent), nothing remains to profile and
both modes coincide -- the scans in \cref{ch:results} are genuine 2-D
likelihood surfaces of the model's full free set, not fixed-nuisance
slices.
\end{theorybox}

\begin{implbox}
\pyfunc{loglik\_grid(model, theta\_baseline, param\_x, x\_grid, param\_y,
y\_grid, value, sigma\_minus=None, sigma\_plus=None, likelihood=...,
profile\_others=False)} returns a real tensor shaped \code{(nx,ny)};
subtracting its minimum and comparing against standard 2-D $\Delta\chi^2$
confidence levels draws confidence contours.
\end{implbox}


% =============================================================================
\chapter{Results}
\label{ch:results}

This chapter presents, interprets, and summarises the real numerical
results of \texttt{notebooks/inference/inference4\_dayabay.ipynb}, executed
end to end against the real Daya Bay 8-detector, 8AD-period data
(\cref{ch:data-release}) using the physics of \cref{ch:detector-implementation}
and the inference machinery of \cref{ch:inference-process}. Every number
below is a real, reproducible output of that notebook, not illustrative.

\section{Real per-reactor relative weight}

\begin{table}[htb]
\centering
\caption{Real per-reactor relative power weight $w_r$ (\labelcref{eq:dayabay-relative-weight}), derived from the real weekly antineutrino-rate history, 8AD period.}
\label{tab:relative-weights-example}
\begin{tabular}{@{}lcccccc@{}}
\toprule
Reactor & R1 & R2 & R3 & R4 & R5 & R6 \\
\midrule
$w_r$ & 1.0010 & 0.9905 & 1.0126 & 0.9965 & 1.0079 & 0.9914 \\
\bottomrule
\end{tabular}
\end{table}

\Cref{tab:relative-weights-example} shows the real values obtained: a
small but genuine $\pm1$--$1.5\%$ real reactor-to-reactor spread around
the 6-reactor average of 1, replacing the previous implicit assumption
that all 6 real cores contribute identically to the flux sum
(\cref{eq:dayabay-face-flux}).

\section{Closure at Daya Bay's own published truth}
\label{sec:results-closure}

\begin{resultbox}
Evaluating the full real 8-detector model at Daya Bay's own published
$\theta_{13}$/$\Delta m^2_{31}$, with $\nu_{\rm norm}=1$ (no fit), gives
\begin{equation}
  N_{\rm pred}^{\rm total} = 3{,}705{,}407, \qquad
  N_{\rm obs}^{\rm total} = 3{,}114{,}934, \qquad
  \frac{N_{\rm pred}}{N_{\rm obs}} = 1.190,
\end{equation}
i.e.\ the simulation over-predicts the real observed rate by
$\approx19\%$, remarkably uniformly across all 8 detectors
($1.179$--$1.195$, \cref{tab:per-detector-ratio-truth}). Gradients
w.r.t.\ $\theta_{13}$, $\Delta m^2_{31}$, and $\nu_{\rm norm}$ are all
finite and non-zero, confirming autograd flows correctly through the real
multi-reactor flux, the real order-1/M cross section, and the real
IAV+LSNL+Gaussian response chain before any fit is attempted.
\end{resultbox}

\begin{table}[htb]
\centering
\caption{Predicted/observed ratio per detector at Daya Bay's own published truth, $\nu_{\rm norm}=1$.}
\label{tab:per-detector-ratio-truth}
\begin{tabular}{@{}lrrr@{}}
\toprule
Detector & Predicted & Real observed & Ratio \\
\midrule
AD11 & 827{,}463  & 695{,}290 & 1.190 \\
AD12 & 840{,}090  & 704{,}771 & 1.192 \\
AD21 & 792{,}729  & 667{,}823 & 1.187 \\
AD22 & 782{,}829  & 658{,}020 & 1.190 \\
AD31 & 116{,}030  &  97{,}069 & 1.195 \\
AD32 & 116{,}150  &  97{,}504 & 1.191 \\
AD33 & 115{,}008  &  96{,}793 & 1.188 \\
AD34 & 115{,}108  &  97{,}664 & 1.179 \\
\bottomrule
\end{tabular}
\end{table}

\section{Absolute-normalization audit}
\label{sec:results-audit}

A dedicated audit of the real $\sim19\%$ gap above checked every component
of the absolute-rate chain: the real target proton count and detector
efficiency (both loaded directly from the official release, not derived);
the order-1/M cross section (independently validated to $\sim1\%$ against
Vogel \& Beacom's own quoted coefficient, \cref{ch:detector-implementation});
the real Huber-Mueller antineutrino yield (the real table spans only
$1.75$--$13.0$~MeV, integrating to $\approx2.0$ antineutrinos/fission above
the IBD threshold -- the correct expected value for that truncated real
range, not a bug); and the thermal-power-to-fission-rate conversion (a
standard textbook formula). No unit or normalization error was found.

\begin{resultbox}
\Cref{tab:table5-crosscheck} cross-checks this notebook's real observed and
predicted rates-per-day against an \emph{independently published} table
(Table~5 of \cite{An2017DayaBayFluxSpectrum}, a different, earlier
621-day dataset). The real \emph{observed} rate matches the published table
to $1$--$3\%$ for every detector, confirming the real data/exposure loading
is correct; the \emph{predicted} rate exceeds it by the same
$\sim19$--$22\%$ already seen above, confirming the gap lives entirely in
the forward-model normalization chain, not in a data-loading or exposure
error.
\end{resultbox}

\begin{table}[htb]
\centering
\caption{Real observed/predicted rate per day vs.\ an independently published rate table \parencite{An2017DayaBayFluxSpectrum}.}
\label{tab:table5-crosscheck}
\begin{tabular}{@{}lrrrrr@{}}
\toprule
Detector & tpeanuts [1/day] & Real obs.\ [1/day] & Published [1/day] & obs/pub. & pred/pub. \\
\midrule
AD11 & 796.0 & 668.9 & 657.18 & 1.018 & 1.211 \\
AD12 & 810.7 & 680.1 & 670.14 & 1.015 & 1.210 \\
AD21 & 723.6 & 609.6 & 594.78 & 1.025 & 1.217 \\
AD22 & 714.9 & 600.9 & 590.81 & 1.017 & 1.210 \\
AD31 &  90.3 &  75.6 &  73.90 & 1.022 & 1.222 \\
AD32 &  90.4 &  75.9 &  74.49 & 1.019 & 1.214 \\
AD33 &  89.5 &  75.4 &  73.58 & 1.024 & 1.217 \\
AD34 &  89.6 &  76.0 &  75.15 & 1.012 & 1.193 \\
\bottomrule
\end{tabular}
\end{table}

Part of the real gap is a known, published physical effect: the
Huber-Mueller flux model itself is measured to predict $\approx5$--$6\%$
more than the real observed rate -- the ``Reactor Antineutrino Anomaly''
\parencite{An2017DayaBayFluxSpectrum}, quoted there as a measured/predicted
flux ratio of $0.946\pm0.020$ for the Huber+Mueller model. The remainder is
consistent with the many further real correlated systematics (target-mass,
efficiency, and flux-normalization nuisances) Daya Bay's own official
combined fit also carries and this project does not reproduce individually
-- precisely why the real official $\nu_{\rm norm}$ nuisance
(\cref{sec:nuisances}) exists rather than being fixed to 1.

\section{Full-spectral-shape fit: three attempts, one honest diagnosis}
\label{sec:results-fullshape}

\begin{figure}[htb]
\centering
\includegraphics[width=0.92\textwidth]{fig1_predicted_vs_real_spectra}
\caption{Predicted (tpeanuts, at Daya Bay's own published truth) vs.\ real
observed prompt-energy spectra, near-hall (AD11) and far-hall (AD31)
detectors.}
\label{fig:predicted-vs-real-spectra}
\end{figure}

\Cref{fig:predicted-vs-real-spectra} shows the real observed spectrum
against the \texttt{tpeanuts} simulation at Daya Bay's own published
truth for one near-hall and one far-hall detector: both follow the same
overall shape (rising then falling with a broad IBD-kinematics-driven peak
around 3--4~MeV), with the far-hall spectrum visibly suppressed relative
to its no-oscillation expectation, the real qualitative signature this
whole measurement is built on.

Minimising \labelcref{eq:poisson-nll} jointly over
$(\theta_{13},\Delta m^2_{31},\nu_{\rm norm})$ against the real
concatenated 8-detector, 208-bin spectrum was attempted three times, with
increasingly rich nuisance-parameter freedom, summarised in
\cref{tab:fullshape-attempts}.

\begin{table}[htb]
\centering
\caption{Full-spectral-shape fit: three real attempts (Section~5.1-5.3 of the notebook).}
\label{tab:fullshape-attempts}
\begin{tabular}{@{}lrrrr@{}}
\toprule
Attempt & $\hat\theta_{13}$ [deg] & $\Delta\hat m^2_{31}$ [$10^{-3}$eV$^2$] & $\hat\nu_{\rm norm}$ & $-2\ln L$ final \\
\midrule
Published truth & 8.506 & 2.528 & -- & -- \\
Unconstrained (5.1) & 9.172 & 7.483 & 0.874 & 6445.8 \\
+ Background nuisances (5.2) & 9.219 & 7.497 & 0.875 & 6390.6 \\
+ LSNL pull nuisances (5.3) & \textbf{0.000} & 11.872 & 0.819 & 4698.4 \\
\bottomrule
\end{tabular}
\end{table}

\begin{resultbox}
\textbf{None of the three attempts converges to a physical point.} The
unconstrained fit (5.1) already diverges: $\hat\theta_{13}=9.17^\circ$ ($+8\%$
absolute, $+15.7\%$ in $\sin^2 2\theta_{13}$) and
$\Delta\hat m^2_{31}=7.48\times10^{-3}$~eV$^2$, roughly $3\times$ the real
published value. Adding the 5 real Gaussian-constrained background-rate
nuisances (5.2, prior widths from \labelcref{eq:background-sigma}) improves
$-2\ln L$ only marginally (6390.6 vs.\ 6445.8, out of a $\sim90{,}000$-unit
total drop from the starting point) and leaves $\hat\theta_{13}$/$\Delta
\hat m^2_{31}$ essentially unchanged -- background-rate uncertainty is not
the dominant residual. Adding the 4 real LSNL pull-curve nuisances instead
(5.3) makes the fit \emph{worse}: $\hat\theta_{13}$ collapses to
$0^\circ$ rather than moving towards $8.51^\circ$, with two of the four
pulls at $3$--$5\sigma$ from their real unit-Gaussian prior. This is a
genuine degeneracy -- LSNL energy-scale freedom can disguise the
oscillation deficit itself as an energy-scale distortion when no
covariance structure constrains it jointly with the oscillation
parameters -- not evidence that LSNL uncertainty does not matter.
\end{resultbox}

\begin{notebox}
This three-way result is itself informative: it rules out both of the two
simplest candidate fixes (background normalization, LSNL curve freedom)
for the full-shape fit and points to the real, still-unimplemented next
step -- a proper correlated covariance or near/far-structured statistical
treatment (\cref{sec:results-summary}) -- rather than more isolated free
parameters.
\end{notebox}

\section{Rate-only fit}
\label{sec:results-rateonly}

\begin{figure}[htb]
\centering
\includegraphics[width=0.75\textwidth]{fig2_rateonly_fit_history}
\caption{Rate-only fit convergence: $-2\ln L$ (Poisson) at every LBFGS closure evaluation.}
\label{fig:rateonly-history}
\end{figure}

Integrating each detector's real spectrum to a single total count
(\labelcref{eq:rate-only-observable}) and fitting $\theta_{13}$/$\nu_{\rm
norm}$ jointly (with $\Delta m^2_{31}$ fixed at Daya Bay's own published
value) against the 8 real total counts:

\begin{resultbox}
\begin{equation}
  \hat\theta_{13} = 8.389^\circ \pm 0.139^\circ, \qquad
  \hat\nu_{\rm norm} = 0.8381\pm0.0008,
\end{equation}
$\sin^2 2\theta_{13}^{\rm fit}=0.08332$ vs.\ Daya Bay's real published
$0.08560$ -- a relative difference of $-2.66\%$, pull $-0.85\sigma$. Every
per-detector predicted/observed ratio at the fitted point lies within
$0.992$--$1.006$ of unity (\cref{fig:rateonly-2dscan}), a real, close
closure -- unlike \cref{sec:results-fullshape}.
\end{resultbox}

\begin{figure}[htb]
\centering
\includegraphics[width=0.62\textwidth]{fig3_rateonly_2d_scan}
\caption{2-D $-2\ln L$ scan over $(\theta_{13},\nu_{\rm norm})$, the rate-only model's full 2-parameter free set.}
\label{fig:rateonly-2dscan}
\end{figure}

\section{Near/far ratio fit: the headline result}
\label{sec:results-nearfar}

Fitting the real far-hall/near-hall per-bin count ratio
(\labelcref{eq:nearfar-ratio}) is the only strategy in this document with
genuine \emph{joint} sensitivity to both $\theta_{13}$ and $\Delta
m^2_{31}$ at once, and it is the only one that converges to a physical
point:

\begin{resultbox}
\begin{align}
  \hat\theta_{13} &= 8.445^\circ \pm 0.134^\circ,
  &\sin^2 2\theta_{13}^{\rm fit}&=0.08442, \\
  \Delta\hat m^2_{31} &= (2.4775\pm0.0548)\times10^{-3}\ {\rm eV}^2. &&
\end{align}
Relative to Daya Bay's own published values: $\sin^2 2\theta_{13}$
differs by $-1.38\%$, $\Delta m^2_{31}$ by $-2.01\%$ -- both within
$\sim1\sigma$. $-2\ln L$ falls from $114.3$ to $26.2$ over a real
26-bin ratio spectrum, a well-conditioned statistic (unlike the
full-shape fit's $\mathcal O(10^5)$ starting value, reflecting the
Poisson error on $\sim3\times10^6$ absolute counts).
\end{resultbox}

\begin{figure}[htb]
\centering
\includegraphics[width=0.85\textwidth]{fig6_nearfar_ratio_comparison}
\caption{Real observed far/near ratio vs.\ \texttt{tpeanuts} simulation (Daya Bay truth) and \texttt{tpeanuts} near/far fit.}
\label{fig:nearfar-ratio-comparison}
\end{figure}

\Cref{fig:nearfar-ratio-comparison} compares the real observed ratio
spectrum against the \texttt{tpeanuts} simulation at Daya Bay's own
published truth (no fit) and the \texttt{tpeanuts} near/far fit: even the
unfit truth-simulation curve already lies close to the real data (unlike
the absolute-count comparison of \cref{sec:results-closure}), because the
$\sim19\%$ normalization gap and much of the residual response mismatch
cancel by construction in the ratio.

\begin{figure}[htb]
\centering
\includegraphics[width=0.62\textwidth]{fig7_nearfar_2d_scan}
\caption{2-D $-2\ln L$ scan over $(\theta_{13},\Delta m^2_{31})$, the near/far model's full 2-parameter free set.}
\label{fig:nearfar-2dscan}
\end{figure}

\Cref{fig:nearfar-2dscan} shows the corresponding 2-D likelihood scan: the
\texttt{tpeanuts} fitted point and Daya Bay's own published point both sit
well inside the same low-$\Delta(-2\ln L)$ region, a direct visual
confirmation of the numerical agreement quoted above.

\begin{figure}[htb]
\centering
\includegraphics[width=0.75\textwidth]{fig5_nearfar_fit_history}
\caption{Near/far ratio fit convergence: $-2\ln L$ (chi2\_asymmetric) at every LBFGS closure evaluation.}
\label{fig:nearfar-history}
\end{figure}

\section{Real data comparison}

\begin{figure}[htb]
\centering
\includegraphics[width=0.85\textwidth]{fig4_real_comparison}
\caption{Real per-detector observed totals vs.\ \texttt{tpeanuts} simulation (Daya Bay truth) and \texttt{tpeanuts} rate-only fit.}
\label{fig:real-comparison}
\end{figure}

\Cref{fig:real-comparison} compares real per-detector observed totals
against the \texttt{tpeanuts} truth-simulation (visibly offset high by the
real $\sim19\%$ normalization gap, \cref{sec:results-audit}) and the
\texttt{tpeanuts} rate-only fit (visibly closer to the real data, having
absorbed that gap into $\hat\nu_{\rm norm}$).

\section{Summary of inference results and comparison with Daya Bay and NuFIT}
\label{sec:results-summary}

\Cref{tab:final-summary} collects every real fitted value obtained in this
chapter alongside Daya Bay's own published best fit
\parencite{An2023DayaBayPrecision} and, for $\theta_{13}$/$\Delta
m^2_{31}$, the independent global-fit value from NuFIT~6.1 (normal
ordering) \parencite{EstebanNuFit2024}, as used elsewhere in this project's
own preset registry (\code{tpeanuts.config.presets.\_SM\_NUFIT61\_NO}).

\begin{table}[htb]
\centering
\caption{Final comparison: tpeanuts inference vs.\ Daya Bay published vs.\ NuFIT 6.1.}
\label{tab:final-summary}
\begin{tabular}{@{}lrrr@{}}
\toprule
Quantity & tpeanuts (near/far, \S\ref{sec:results-nearfar}) & Daya Bay published & NuFIT 6.1$^\dagger$ \\
\midrule
$\sin^2 2\theta_{13}$ & $0.08442$ & $0.0856$ \parencite{An2023DayaBayPrecision} & $0.0878$ \\
$\theta_{13}$ [deg] & $8.445\pm0.134$ & $8.506$ & $8.62$ \\
$\Delta m^2_{31}$ [$10^{-3}$eV$^2$] & $2.478\pm0.055$ & $2.528$ & $2.511$ \\
\bottomrule
\end{tabular}
\vspace{2pt}

{\footnotesize $^\dagger$NuFIT 6.1's own global-fit best-fit central value
(normal ordering; itself informed by, among other inputs, Daya Bay's real
published measurement, not an independent cross-check) \parencite{EstebanNuFit2024}.}
\end{table}

\begin{resultbox}
\textbf{Final assessment.} The real near/far ratio fit
(\cref{sec:results-nearfar}) recovers both $\theta_{13}$ and $\Delta
m^2_{31}$ jointly, within $1$--$2\%$ of Daya Bay's own published values and
comfortably within NuFIT's global-fit range -- using \texttt{tpeanuts}'s
own real 6-reactor flux, real order-1/M cross section, real detector
response, and Daya Bay's own real near/far discovery strategy
\parencite{An2012DayaBayDiscovery}, not a synthetic or illustrative
dataset. The rate-only fit independently confirms $\theta_{13}$ alone to
similar precision. The full-spectral-shape fit, despite three real
attempts (including two different real nuisance-parameter families), does
not converge to a physical point, and this is reported as a genuine,
diagnosed finding rather than hidden: it demonstrates precisely why real
reactor-antineutrino collaborations build their combined analyses around a
full correlated systematic covariance, rather than a handful of
independent nuisance parameters, once the statistics are large enough that
sub-percent shape residuals dominate over the Poisson error.

\textbf{Recommended next step.} A real near/far \emph{spectral-shape}
covariance treatment (correlating the response/flux systematics between
the near and far halls directly, rather than only comparing their bin-wise
ratio) is the most promising avenue to make the full-spectral-shape fit
itself robust, now that both simpler nuisance-parameter fixes have been
tried and shown insufficient.
\end{resultbox}


\clearpage
\printbibliography[title={References}]

\end{document}
